\section{Local experimentation and the decision to use HPC}
\label{architecture-local-experiments}

A central goal is high-fidelity atmospheric simulation on ordinary workstations
and MacBooks through efficient algorithms, careful memory ownership and GPU
acceleration. Local work is not intended to stop at toy models, coarse previews
or preliminary checks before an inevitable cluster run. Reducing avoidable
work and data movement can make substantially more demanding experiments
practical on the same hardware. A researcher
should be able to change a documented experiment, inspect its assumptions,
and increase its resolution without maintaining a separate solver fork for
each case. The objective is useful scientific work per unit of time, memory
and effort. A high GPU utilization percentage is not sufficient evidence of
that objective, and a run does not become scientifically adequate merely
because it completes on a laptop or a cluster.

Three questions determine where an experiment belongs. First, does the
discretization and selected physics answer the scientific question? Second,
does the complete working set fit the available resources? Third, can the
required duration, repetitions and analysis finish within the researcher's
time budget? A configuration can pass one question and fail another. The
\href{MATHEMATICS.pdf}{mathematics guide's resolution and experiment-design
chapter} develops the first question and the scaling assumptions behind the
other two. This section describes the engineering boundaries that keep the
answers separate.

\subsection{Capacity and turnaround are different constraints}

A useful memory estimate follows simultaneously live owners. It includes
primary fields, time-integration stages, scheme-private state, geometry,
rollback candidates, reports, output snapshots and communication buffers.
The simulator, Studio, Python and operating system can all be active at once.
On a discrete-memory machine, RAM and each GPU's local memory are separate
constraints. On a unified-memory machine, competing allocations can consume
one physical capacity even though several subsystems each report a plausible
local budget. Neither topology removes maximum allocation sizes or the need
for headroom.

For illustration, a float32 scalar array of size $512^3$ occupies 0.5 GiB.
The existing 51-field-equivalent dynamics admission envelope would account
for 25.5 GiB for that allocation shape before interpreting additional owners
and host/device duplication. Doubling all three dimensions makes a scalar
array 4 GiB and that envelope 204 GiB. These are arithmetic examples, not
recommended Nimbus grids, measured peak usage, or predictions that a device
with that nominal capacity will admit the job. The
\href{../include/compute/compute_memory_admission.hpp}{memory admission owner}
must also check the largest buffer, backend capabilities and simultaneous
levels. A preflight estimate should explain which owners it covers; it cannot
serve as a measurement of peak process memory.

Turnaround includes more than stepping. A useful comparison is
\[
 T_{\mathrm{local}}=T_{\mathrm{prepare}}+T_{\mathrm{solve}}
                      +T_{\mathrm{publish}}+T_{\mathrm{analyse}},
\]
with overlapping phases counted on the actual critical path rather than
added twice. Remote execution also introduces queueing, staging, result
transfer and recovery costs. A cluster can shorten the solve and still delay
the answer to a small experiment. Conversely, an overnight local run may be
unsuitable when the scientific conclusion needs dozens of realizations.
Measure the whole workflow for the same experiment before choosing placement.

\subsection{Choose the kind of parallelism that serves the question}

\begin{longtable}{@{}>{\raggedright\arraybackslash}p{0.22\textwidth}>{\raggedright\arraybackslash}p{0.34\textwidth}>{\raggedright\arraybackslash}p{0.36\textwidth}@{}}
\toprule
Need & Useful execution arrangement & Consequence \\
\midrule
\endhead
One admitted case fits locally & CPU or one selected accelerator & Fast
iteration can avoid remote staging; actual elapsed time still needs measurement. \\
Many independent cases fit individually & Separate processes or batch jobs &
Distribute experiments without halo exchange; each job still needs a complete
working set and a distinct output identity. \\
One case exceeds a device's memory & Spatial domain decomposition & Requires
the distributed ownership and exchange design in
\hyperref[architecture-multigpu]{the multi-GPU section}; adding cards alone
does not implement it. \\
One case fits but misses the time budget & Profile first, then consider
strong scaling & Transfers, reductions and output can erase a faster kernel
or the benefit of an additional device. \\
The ensemble or archive exceeds local capacity & Managed remote resources &
Scheduling, storage, transfer and reproducible restart become part of the
experiment, not administrative details. \\
\bottomrule
\end{longtable}

Running independent cases is often the first useful step beyond one machine.
They can share immutable coefficient files but must not share writable run
directories, mutable model state or a GPU allocation budget without explicit
coordination. Device assignment and the number of simultaneous jobs must be
controlled. The current remote-compute workflow launches a process through
SSH; it does not provide a distributed solver or a general cluster scheduler.
An ensemble diagram in Studio is an analysis facility, not evidence that a
multi-node ensemble scheduler exists.

Local alternatives include a shorter diagnostic experiment, a different output
cadence, a smaller scientifically justified domain, or an admitted nest focused
on the region of interest. Each changes a different cost. A shorter duration
cannot answer a question about late-time behavior. Smaller domains change
boundary influence. A coarser grid changes resolved scales. Nesting introduces
exchange and feedback requirements and does not bypass their readiness gates.
These choices must remain visible in the experiment record rather than being
silently selected by an automatic speed policy.

\subsection{What optimization can buy without changing the experiment}

Memory layout, traversal order and cache reuse can reduce the cost of moving
the same physical state. Reusing valid material coefficients can avoid
recomputing them at every acoustic operation. Keeping intermediate state on
the accelerator can avoid unnecessary host round trips. Reusing scratch
capacity can reduce allocation overhead, while shortening actual buffer
lifetimes can reduce peak memory. These improvements answer different
bottlenecks and can combine to make a demanding local experiment practical.
They must retain the producer/consumer, stage and failure contracts developed
in the execution chapters.

GPU acceleration is most useful when enough suitable work can remain on the
device and the selected numerical families have native implementations.
Launching many tiny kernels or repeatedly copying complete fields through
the host can consume the anticipated gain. Conversely, a small diagnostic or
control task can be inexpensive on the CPU; moving every function to a shader
is not itself the objective. Measure which owner limits the complete run,
then remove that cost without silently changing the selected physics, output
or acceptance conditions.

The result can be invested in finer resolution, longer duration, more ensemble
members or faster interactive iteration. The mathematics guide explains why
linear-resolution gains are nonlinear in runtime, and why memory and time
must both be considered. A faster implementation and a smaller live working
set can shift the practical local/HPC boundary. That boundary should be
re-measured after optimization rather than inferred from the resource demands
of an early implementation.

\subsection{Changing an experiment without rewriting its infrastructure}

Existing public choices should be expressed through the versioned configuration
and recorded with the executable catalog. Initial conditions, physical domain,
boundary conditions, admitted physics, output and backend selection belong to
that document. A case-specific condition hidden in an inner loop cannot be
recovered by reading the saved YAML and is therefore a poor extension boundary.

A new physical formulation still requires implementation. The reusable unit
is its state and operator contract: the quantities and units it reads, the
inventories it changes, grid locations, source interval, private continuation
state, invalid-input behavior and supported backends. It must join ordinary
configuration loading, admission, restart, diagnostics and tests. If it requires
native execution, a host implementation alone does not complete the feature.
The existing \href{../include/compute/backend_interface.hpp}{typed backend
interface} and \href{../src/core/runtime/scheme_catalog.cpp}{scheme catalog}
provide boundaries to extend; they are not a dynamic plugin system that can
accept arbitrary new physics without compiled code.

This also separates analysis customization from solver customization. A new
Python diagnostic can operate on published fields without changing the time
integrator. It must disclose missing inputs, units, grid reconstruction and
sampling. A new prognostic variable has a much larger obligation: initialization,
transport, source coupling, halos, restart and publication all need an owner.
Modern engineering reduces the repeated plumbing, but does not make those
scientific obligations optional.

\subsection{Measure a decision boundary, not a universal machine ranking}

A useful future resource report would combine observed peak working sets,
stage times, transfers, output volume and actual simulated duration with the
configuration identity. It should state uncertainty and identify extrapolation.
The existing admission checks, telemetry and Studio loading-worker measurement
cover pieces of this problem; they do not yet form a universal experiment-cost
estimator. In particular, Studio's small cached NPY measurement chooses a
loading heuristic and cannot choose solver resolution or predict a storm's
runtime.

The practical development order is to establish correctness for a declared
profile, optimize and measure complete single-device work on representative
unified and discrete-memory systems, then evaluate independent-job throughput and finally
coupled distributed execution where the experiment requires it. Lowering the
barrier means keeping the local build usable throughout that process. Optional
cluster transports should not become prerequisites for running a single CPU
case, and a proposed scaling feature should not be advertised as available
before it has public configuration, working coupling and the required tests.
